% ============================================================
% 03_marco_teorico.tex
% ============================================================

\section{Marco Teórico y Antecedentes}
\label{sec:marco}

\subsection{Internet de las Cosas (IoT) en entornos educativos}

El Internet de las Cosas (IoT, por sus siglas en inglés) se define como la red de
objetos físicos equipados con sensores, software y conectividad que les permite
intercambiar datos con otros sistemas y dispositivos \parencite{Atzori2010}. En el
ámbito de los edificios educativos, las aplicaciones IoT se han orientado
principalmente a la gestión energética, el monitoreo de calidad del aire, el control
de acceso y la automatización de la iluminación \parencite{Jia2019}.

\textcite{Shaikh2014} señalan que los sistemas de control de iluminación basados
en sensores de ocupación y luminosidad constituyen uno de los casos de uso IoT con
mayor retorno de inversión en entornos de uso esporádico como aulas universitarias.
La combinación de sensores fotoeléctricos y detectores de presencia permite aplicar
estrategias de control adaptativo que atienden simultáneamente al confort visual y
al ahorro energético.

\subsection{Protocolo MQTT}

MQTT es un protocolo de mensajería ligero basado en el modelo publicador/suscriptor
(\textit{publish/subscribe}). Un elemento central del protocolo es el
\textit{broker}, servidor que gestiona la distribución de mensajes entre clientes
según los temas (\textit{topics}) a los que están suscritos. MQTT opera sobre TCP/IP
y define tres niveles de calidad de servicio (QoS): entrega sin garantía (QoS\,0),
entrega al menos una vez (QoS\,1) y entrega exactamente una vez (QoS\,2)
\parencite{Banks2019}.

La versión 3.1.1 del estándar, publicada por OASIS, es la más ampliamente soportada
en hardware embebido. El cliente MQTT integrado en el ESP-IDF de Espressif
implementa dicha versión e incluye soporte para reconexión automática,
autenticación por usuario y contraseña, y publicaciones con retención
\parencite{Espressif2024}.

\subsection{ESP32-C3 y ESP-IDF}

El ESP32-C3 es un SoC (System-on-Chip) de Espressif Systems basado en la
arquitectura RISC-V de 32 bits, con Wi-Fi 802.11\,b/g/n y Bluetooth 5 (LE)
integrados. Opera a frecuencias de hasta 160\,MHz y dispone de 400\,KB de SRAM y
4\,MB de Flash \parencite{EspressifDS2024}. Para aplicaciones IoT de bajo consumo,
el chip implementa modos \textit{modem sleep}, \textit{light sleep} y
\textit{deep sleep}, con consumos de tan solo 5\,µA en el modo más profundo.

El ESP-IDF (\textit{Espressif IoT Development Framework}) es el framework oficial de
desarrollo para la familia ESP32. Proporciona una capa de abstracción de hardware
(HAL) sobre FreeRTOS, drivers certificados para periféricos —incluyendo I\textsuperscript{2}C
maestro, UART y Wi-Fi—, y un sistema de gestión de componentes que facilita la
incorporación de bibliotecas externas como cJSON y el cliente MQTT
\parencite{Espressif2024}.

\subsection{Sensores utilizados}

\subsubsection{BH1750 — Sensor de luminancia}

El BH1750FVI es un sensor digital de luminancia ambiental de 16 bits fabricado por
Rohm Semiconductor, con interfaz I\textsuperscript{2}C. Opera en el rango de 1 a
65535\,lx con resolución de 1\,lx en modo de alta resolución continua
(\textit{Continuous H-Resolution Mode}). Admite dos direcciones I\textsuperscript{2}C
(0x23 y 0x5C) según el estado del pin ADDR, lo que posibilita la conexión de dos
unidades en el mismo bus \parencite{Rohm2011}.

La conversión analógica-digital de la señal fotoeléctrica interna entrega un valor
crudo de 16 bits que se convierte a lux mediante la ecuación:

\begin{equation}
  E_v = \frac{V_{\text{raw}}}{1{,}2}
  \label{eq:bh1750}
\end{equation}

\noindent donde $E_v$ es la iluminancia en lux y $V_{\text{raw}}$ es el valor
entero de 16 bits leído del sensor.

\subsubsection{SHT30/SHT31 — Sensor de temperatura y humedad}

El SHT30 de Sensirion es un sensor digital de temperatura y humedad relativa con
interfaz I\textsuperscript{2}C. Ofrece precisión de ±0,3\,°C en temperatura
(rango $-$40\,°C a +125\,°C) y ±3\,\%\,HR en humedad \parencite{Sensirion2023}.
El protocolo de comunicación contempla un comando de disparo de medición de 16 bits
y la lectura de 6 bytes de respuesta que incluyen temperatura, humedad y bytes de
verificación CRC-8.

Las fórmulas de conversión de los valores crudos de 16 bits son:

\begin{equation}
  T\,[\mathrm{{}^{\circ}C}] = -45{,}0 + 175{,}0 \cdot \frac{V_T}{2^{16}-1}
  \label{eq:sht_temp}
\end{equation}

\begin{equation}
  HR\,[\%] = 100{,}0 \cdot \frac{V_H}{2^{16}-1}
  \label{eq:sht_hum}
\end{equation}

\noindent donde $V_T$ y $V_H$ son los valores crudos de temperatura y humedad,
respectivamente \parencite{Sensirion2023}.

\subsubsection{HLK-LD2420 — Radar de presencia}

El HLK-LD2420 de Hi-Link es un módulo de radar FMCW (\textit{Frequency-Modulated
Continuous Wave}) de 24\,GHz diseñado para detección de presencia humana. A
diferencia de los sensores PIR (infrarrojo pasivo), el radar FMCW detecta tanto
movimiento macroscópico como la microvibración respiratoria de personas estáticas
\parencite{HiLink2023}. Opera en el rango de 0 a 8\,m con resolución de 0,7\,m y
se comunica mediante UART a 115200 baudios.

En modo de reporte nativo simple (\textit{Native Simple Report Mode}), el módulo
emite tramas de texto con tres tipos de mensajes: \texttt{ON} (presencia detectada),
\texttt{OFF} (sin presencia) y \texttt{Range }\textit{d} (distancia al objeto en
centímetros). Este modo reduce la complejidad del procesamiento en el
microcontrolador respecto al modo de protocolo binario.

\subsection{Node-RED}

Node-RED es una herramienta de programación visual basada en flujos para aplicaciones
IoT, desarrollada originalmente por IBM y actualmente mantenida por la OpenJS
Foundation \parencite{NodeRED2024}. Permite construir aplicaciones conectando nodos
funcionales de forma gráfica, con soporte nativo para protocolos MQTT, HTTP y WebSocket.
Implementado sobre Node.js, su ejecución en Raspberry Pi dentro de contenedores Docker
simplifica el despliegue y la portabilidad del sistema.

\subsection{Antecedentes relevantes}

\textcite{Cardiel2012} propusieron una arquitectura de middleware orientada a SCADA
basada en el modelo de publicación/suscripción, conceptualmente similar a la empleada
en este trabajo, con separación entre la capa de adquisición de datos y la capa de
lógica de negocio. \textcite{Jia2019} describieron un sistema IoT para gestión
energética de edificios educativos usando sensores de ocupación e iluminación,
obteniendo reducciones de consumo superiores al 40\,\%. Estos antecedentes respaldan
la viabilidad del enfoque adoptado en NodeAlert IoT.
